%!TeX root=../main.tex

\chapter{Metodologia}
\label{cap:metodologia}

Para avaliar o sistema proposto neste trabalho, foi realizado um experimento com 15 participantes sem conhecimento prévio de shogi, durante um período de três semanas. Ao longo desse período, os participantes foram instruídos a resolver pelo menos uma lista de exercícios por dia, utilizando o sistema.

Durante o experimento, foram coletadas métricas relacionadas ao desempenho dos participantes, incluindo o tempo de resposta, a taxa média de acertos por módulo e a quantidade de listas necessárias para avançar para o módulo seguinte. Essas métricas permitiram analisar o desempenho dos usuários e o funcionamento da progressão de dificuldade implementada no sistema.

Ao final do experimento, foi disponibilizado um questionário para coletar as opiniões dos participantes sobre diferentes aspectos do sistema, como usabilidade, organização dos módulos, dificuldade dos exercícios e contribuição para o aprendizado. O questionário foi composto pelas seguintes perguntas:

\begin{itemize}
    \item De 1 a 10, você considera a interface intuitiva e fácil de utilizar?
    \item Qual foi o módulo que você teve mais dificuldade?
    \item Qual o motivo?
    \item A progressão dos módulos faz sentido? Por que?
    \item Você considera que a quantidade de exercícios por lista é a ideal? Por que?
    \item Você considera que aprendeu os fundamentos básicos do shogi após utilizar o sistema? Por que?
    \item De 1 a 10, você considera que as explicações apresentadas pelo sistema são claras e fáceis de compreender?
    \item De 1 a 10, o quanto você considera que as explicações ao final das listas contribuiu para o seu aprendizado?
    \item Você tem alguma sugestão de pontos a melhorar no sistema?
\end{itemize}